\section{Introduction}


Deploying large language models (LLMs) directly on edge devices is crucial. On-device usage eliminates delays caused by sending data to a cloud server and enables LLMs to operate offline, which is beneficial for real-time applications like virtual assistants, chatbots, and autonomous vehicles. The operational costs associated with maintaining and scaling centralized cloud infrastructure can also be reduced. On-device LLM also enhances data security by keeping sensitive information local, reducing the chance of data breaches. LLMs, grounded in transformer-based architectures~\cite{vaswani2017attention}, have gathered significant attention for their impressive performance across diverse benchmarks~\cite{gpt3, opt, touvron2023llama, scao2022bloom}. However, the large model size leads to the high serving costs. For example, GPT-3 has 175B parameters, which is 350GB in FP16, while the latest B200 GPU only has 192GB memory, let alone edge devices.

\begin{figure}
    \centering
     \includegraphics[width=0.5\textwidth]{figures/teaser.pdf}
    \caption{We introduce \textbf{AWQ}, a versatile weight quantization method for LLM. To implement AWQ, we developed \textbf{\system} to deploy 4-bit quantized LLMs into various edge platforms, achieving a \textbf{3-4$\times$} performance boost compared to FP16. Notably, we've also manufactured a \textbf{TinyChat computer}, powered by \system, which contains an NVIDIA Jetson Orin Nano with only 8GB of memory and 15W power consumption. Demo: \url{https://youtu.be/z91a8DrfgEw}.} 
    \label{fig:teaser}
\end{figure}


Low-bit weight quantization for LLMs can significantly reduce the memory footprint of on-device LLM inference but is hard.
Quantization-aware training (QAT) is not efficient due to the high training cost, while post-training quantization (PTQ) suffers from large accuracy degradation under a low-bit setting. The closest work is GPTQ~\cite{frantar2022gptq}, which uses second-order information to perform error compensation. 
However, it may overfit the calibration set during reconstruction, distorting the learned features on out-of-distribution domains (Figure~\ref{fig:calib_set_ablation}), which is problematic since LLMs are \emph{generalist} models. 

In this paper, we propose  \method (\methodshort), a hardware-friendly low-bit weight-only quantization method for LLMs. Our method is based on the observation that \emph{weights are not equally important} for LLMs' performance. There is a small fraction (0.1\%-1\%) of \emph{salient} weights; skipping the quantization of these salient weights will significantly reduce the quantization loss (Table~\ref{tab:fp_ratio}). To find the salient weight channels, the insight is that we should refer to the \emph{activation} distribution instead of the \emph{weight} distribution, despite we are doing \emph{weight-only} quantization: weight channels corresponding to larger activation magnitudes are more salient since they process more important features. 
To avoid the hardware-inefficient mixed-precision implementation, we analyze the error from weight quantization and derive that \emph{scaling up the salient channels can reduce their relative quantization error} (Equation~\ref{eq:scaled_quant}). Following the intuition, we designed a per-channel scaling method to automatically search for the optimal scaling that minimizes the quantization error under full-weight quantization.
\methodshort does not rely on any backpropagation or reconstruction, so it can well preserve LLMs’ generalization ability on various domains and modalities without overfitting to the calibration set. 

To implement AWQ, we designed \system, an efficient inference framework to convert theoretical memory savings from 4-bit LLM to measured speedup. Our framework significantly speeds up linear layers through on-the-fly dequantization. We also take advantage of efficient 4-bit weight packing and kernel fusion to minimize the inference overhead (\eg, intermediate DRAM access and kernel launch overhead), such that we can better realize the speed up from quantizing the weights to 4-bit, despite the computer is byte-aligned. 


Experiments show that \methodshort outperforms existing work on various tasks
for  different model families (\eg, LLaMA~\cite{touvron2023llama}, OPT~\cite{opt}) and model sizes.  Thanks to better generalization, it also achieves good quantization performance for \emph{instruction-tuned} LMs (\eg, Vicuna) and, for the first time, \emph{multi-modal} LMs (OpenFlamingo~\cite{openflamingo}). 
\system further translates the $\sim$4$\times$ lower memory footprint to measured speedup. On desktop, laptop and mobile GPUs, we consistently observe a \textbf{3.2-3.3}$\times$ average speedup compared to the FP16 implementation by Huggingface across a diverse spectrum of LLMs. Furthermore, it facilitates effortless deployment of the Llama-2-70B model on a single NVIDIA Jetson Orin with 64GB of memory. It also democratizes 13 billion parameter LLM at an interactive pace of 30 tokens/second on a laptop RTX 4070 GPU with only 8GB of memory.  \methodshort has been widely adopted by industry and open-source community: \href{https://huggingface.co/docs/transformers/main_classes/quantization}{HuggingFace Transformers}, \href{https://github.com/NVIDIA/TensorRT-LLM/}{NVIDIA TensorRT-LLM}, \href{https://blogs.windows.com/windowsdeveloper/2024/05/24/quantization-with-directml-helps-you-scale-further-on-windows/}{Microsfot DirectML}, \href{https://console.cloud.google.com/vertex-ai/publishers/meta/model-garden/llama-2-quantized}{Google Vertex AI}, \href{https://github.com/intel/neural-compressor}{Intel Neural Compressor}, \href{https://aws.amazon.com/blogs/machine-learning/boost-inference-performance-for-llms-with-new-amazon-sagemaker-containers/}{Amazon Sagemaker}, \href{https://community.amd.com/t5/ai/reduce-memory-footprint-and-improve-performance-running-llms-on/ba-p/686157}{AMD}, \href{https://github.com/lm-sys/FastChat/blob/main/docs/awq.md}{FastChat}, \href{https://github.com/vllm-project/vllm/blob/main/vllm/model_executor/layers/quantization/awq.py}{vLLM}, \href{https://github.com/InternLM/lmdeploy}{LMDeploy}, and enables Falcon-180B deployable on a \href{https://github.com/NVIDIA/TensorRT-LLM/blob/main/docs/source/blogs/Falcon180B-H200.md}{single} H200 GPU.






